\begin{refsection}
\chapter{QR factorisation}

This chapter describes an implementation of QR factorisation and related
algorithms by Oak Leaf.
The implementation is based on ideas, the spirit, and
a framework developed in the book \fullcite{golub2013matrix}
(henceforth cited as \cite{golub2013matrix}).
The fundamental principles of QR decomposition, the algorithm
by John~G.F.~Francis\footfullcite{wiki:jfrancis},
can be found in references, see Section~\ref{sec:qrbiblio}.

I have been studying QR factorisation in the hope I will find
genuinely deep understanding of principles, and properties, of the approach.
It has helped me to select of effective algorithms for optimisation problems
solved in Oak Leaf
rather than to apply of general forms of the algorithms.
Another benefit is the implementation in modern Fortran
equipped by vectorisation computing capabilities;
that permits writing of a clean and computationally efficient code.
By the way, I had been able to drop dependency of Oak Leaf on external libraries,
however compilers itselfs utilises routines by Lapack
(\footfullcite{website:lapack}), the well-known project on base
of \cite{golub2013matrix}.

I had developed both computation effective and debug variants of the routines.
A source code for the effective way can be found
in \texttt{src/qrhouse.f08}, whilst \texttt{src/qrdebug.f08}
is the basic, naive, implementation intended for verification of algorithms
at all. Run-time tests contains the files \texttt{test/qr/testqrhouse.f08},
and \texttt{test/qr/testqrdebug.f08}.

\section{QR factorisation}
\label{sec:qrfactorisation}

QR algorithm factorise (decompose) of a matrix $A$ on:
an upper diagonal triangular matrix $R$, and $Q$ matrix
having orthogonal columns
\begin{equation}
  A = Q R.
\end{equation}
The decomposition is a wise way how to solve systems of linear equations,
to derive an inversion, or to find both eigenvalues and eigenvectors
of a matrix. QR factorisation can be applied on non-square matrices,
or rank deficient systems of linear equations;
the problems so much important for real world applications.
All these astonishing abilities raises from the extraordinary coincidence:
a linear system with triangular matrix $R$ can be solved trivialy, and
the inversion of an orthogonal matrix is the matrix transpose $Q^{-1} = Q^T$.

QR factorisation can be applied on the matrix $A^{m \times n}$ having real
elements $a_{i,j} \in \mathcal{R}$ arranged in $m$ rows and $n$ columns.
$i$ is an index for rows, $j$ for columns.
The result matrices has dimensions $m \times n$ for $R$, and
$m \times m$ for $Q$.
$A$ can be non-square $m \ne n$, singular, and asymmetric.
$Q$ is a matrix having orthogonal columns
\begin{align*}
  Q_i^T \cdot Q_j = 0, \quad i \ne j,\\
  Q_i^T \cdot Q_i \ne 0, \quad i = j,
\end{align*}
or $Q^T Q = Q Q^T = I$.
The generalisation on complex elements of $A$ is straightforward,
but out of scope of our interest.

QR factorisation for elements of $3 \times 3$ matrix can be written as
\begin{equation}
    \begin{pmatrix*}[r]
      a_{11}, & a_{12}, & a_{13},\\
      a_{21}, & a_{22}, & a_{23},\\
      a_{31}, & a_{32}, & a_{33}
    \end{pmatrix*}
    =
    \begin{pmatrix*}[r]
      q_{11}, & q_{12}, & q_{13},\\
      q_{21}, & q_{22}, & q_{23},\\
      q_{31}, & q_{32}, & q_{33}
    \end{pmatrix*}
    \begin{pmatrix*}[r]
      r_{11}, & r_{12}, & r_{13},\\
      0,     & r_{22}, & r_{23},\\
      0,     & 0,     & r_{33}
    \end{pmatrix*}.
\end{equation}

The principle of QR algorithm is to find such equivalence transformation
of $A$
which will nullify $r_{ij} = 0$ under-diagonal elements $i < j$ of $R$
together with keeping $Q$ orthogonal.
Householder reflection, represented by
the symmetric $m \times m$ projection matrix
\begin{equation} \label{eq:housemat}
  P = I - \beta \, v \otimes v^T, \quad
  \beta = \frac{2}{v^T \cdot v},
\end{equation}
satisfies the requirement.\footnote{
  The outer and the inner (scalar) product of vectors
  $u,v \in \mathrm{R}^m$ are defined as
  \[
  u \otimes v^T = u_i v_j, \; i,j = 1, \ldots, m,
  \quad \text{and} \quad
  u^T \! \cdot v = \sum_{i=1}^m u_i v_i.
  \]
}
The determination of Householder vector $v$, the heart of the factorisation,
will be described in Section~\ref{sec:house}.

Since this point, Fortran notation for element block selection of matrices
and vectors will be used:
the plain $(:)$ means whole row or column,
$\mdx{j}{m}$ is a part with starting index $j$
and the last index $m$ (including the interval points).


A simple choice for the annihilation of all sub-diagonal elements of
the column $A_j$ is to get $j$-th column of $A$:
$x = A_j \in \mathcal{R}^m$, and to set the Householder vector as
\begin{equation}
  v(\mdx{j}{m}) = x(\mdx{j}{m}) \pm \lVert x(\mdx{j}{m}) \rVert_2 \ e_j
\end{equation}
for the elements of $x$ with indexes $j \ldots n$, zero otherwise.
The vector is a reflection around the unit vector $e_j$.
Householder matrix (\ref{eq:housemat}) for $j=2$ is derived
from $v^{(2)} = (0,v_2,v_3)^T$ in the form
\begin{equation}
  P^{(2)} =
    \begin{pmatrix*}[r]
      1, & 0, & 0,\\
      0, & p_{22}, & p_{23},\\
      0, & p_{32}, & p_{33}
    \end{pmatrix*}.
\end{equation}

The annihilation process is designed by this way: the matrix multiplication
$P^{(1)} A$ eliminates the first sub-column giving
the first iteration of $R^{(1)}$, the updated matrix is used to
the eliminate next column $R^{(2)} = P^{(2)} R^{(1)}$, and the machinery
is repeated until finish in $n$ steps.
\begin{equation}
  R^{(j)} = P^{(j)} R^{(j-1)}, \quad j = 1 \dots n.
\end{equation}
The sequence has start from the full matrix $R^{(0)} = A$,

The orthogonal matrix $Q$ collects the projection matrices
during the process.
\begin{equation}
  Q^{(j)} = Q^{(j)} P^{(j-1)}, \quad j = 1 \dots n
\end{equation}
The sequence has start from a unitary matrix $Q^{(0)} = I$ of dimension $m$.

Algorithm~\ref{alg:qrfactorisation} summarises our development.
One can be implemented via matrix operations.


\begin{algorithm}
  \caption{QR factorisation}
  \label{alg:qrfactorisation}
  \begin{algorithmic}
    \State \Comment{Let $A, R \in \mathcal{R}^{m \times n},
    Q,I,P \in \mathcal{R}^{m \times m}, v \in \mathcal{R}^m, \beta \in \mathcal{R}$}
    \State $R = A$
    \State $Q = I$
    \For {$\; j = 1 \ldots n$}
    \State $v, \beta$ = \Call{housevec}{$R_j$}
    \State $P = I - \beta \, v \otimes v^T$
    \State $R = P R$
    \State $Q = Q P$
    \EndFor
    \State \Return factored $A$ into $R,Q$
  \end{algorithmic}
\end{algorithm}

The $R$ can be computed using of the sub-matrix
$R(\mdx{j}{m},\mdx{j}{n}) = P(\mdx{j}{m},\mdx{j}{m}),R(\mdx{j}{m},\mdx{j}{n})$;
this way of acceleration is not valid for $Q$.

There are more ways how to derive QR factorisation.
Householder projection is the way having best numerical
stability and properties. Givens rotations, or modified
Gram-Schmidt process, are another methods
useful in particular problems.


\section{An effective QR factorisation}

The formulation of QR factorisation in previous section~\ref{sec:qrfactorisation}
can be improved in terms of both speed up and memory economy.

\subsection{Speed up}
\label{sec:speedup}

The computations of $P$ in (\ref{eq:housemat}) are expensive due
the outer product of
Householder vector $v \in \mathcal{R}^m$, followed by
the additional matrix multiplication.

A way how to significantly improve speed of the approach comes
from the observation (\cite{golub2013matrix}, Sect.~5.1.4)
that the projection matrix $P$ is commonly used as part of
matrix products. In such case, $P$ can be eliminated
in behalf of $v$:
\begin{equation} \label{eq:golub514a}
  PA = (I - \beta \, v \otimes v^T) \, A = A - (\beta v) \otimes (v^T \! A)
\end{equation}
as well as
\begin{equation} \label{eq:golub514b}
  A P = A \, (I - \beta \, v \otimes v^T) = A - (A v) \otimes (\beta v)^T.
\end{equation}
The re-arrange removes one step in the above algorithm.
Another step reduction is possible, when $Q$ is not required explicitly.

\subsection{Room sharing}
\label{subsec:memshare}

The second improvement evolves the previous idea:
if we need only $v$ for successive operations, we can drop
to keep matrix $Q$ which can be recovered from vectors $v$,
if need.
That approach is very handy;
otherwise, we should keep both $Q$ and a set of $v$ in order
to be able do any operations with factored matrices.

Now, we will keep only vectors $v$, and recover $Q$ later.
The vectors $v$ can be stored in an additional matrix $(v^{(j)})$,
but there is a better way. All the vectors $v^{(j)}$ has
the structure
\[
(0, \ldots, v_j, v_{j+1}, \ldots, v_m )^T.
\]
with $j-1$ null elements.
The determination of $P$ (\ref{eq:housemat}) is independent
on norm of $v$, so it becomes useful to introduce the normalisation
\begin{equation}
  v = \left( 0, \ldots, 1, \frac{v_{j+1}}{v_j}, \ldots, \frac{v_m}{v_j} \right)^T.
\end{equation}
Only the elements $j+1, \ldots, n$ carries desired information, and we
can use under diagonal elements, the unfilled place, of matrix $R$ to store
parts of the vectors.

The final structure of matrix $R$, which can share memory with
$A$ being overwritten sequentially, will be
\begin{equation}
    \begin{pmatrix*}[r]
      r_{11}, & r_{12}, & r_{13}, & r_{14}, \\
      v_2^{(1)}, & r_{22}, & r_{23}, & r_{24},\\
      v_3^{(1)}, & v_3^{(2)}, & r_{33}, & r_{34},\\
      v_4^{(1)}, & v_4^{(2)}, & v_4^{(3)}, & r_{44}
    \end{pmatrix*}.
\end{equation}


\subsection{Some notes}

It may be disputable, if this way to save some memory is rational
today, but the related speed up should be important for applications.

The algorithm here is not intended for solution of very large
problems with sparse matrices. The expected matrix should be dense
and small. For that matrices, the Householder transformation
offers numerical stability --- it is the most stable known method.

The execution speed can be improved by a parallel implementation.
Unfortunately, the Householder matrix should be applied in an exact
sequence, the parallelization is difficult but possible via
a recursion (see \fullcite{golub2013matrix}, henceforth cited as \cite{golub2013matrix} Sect.~5.2.4).
The more parallel friendly way is to replace
Householder transformations by Givens rotations.
Unfortunately, they brings worse numerical stability so they
are recommended for large sparse problems only.

The effective way, described in this section, can speed-up
computations, but it will work only for compiled computer languages
like Fortran, Go, or Julia.
The scripting languages, Matlab or Python,
will benefits from the matrix formulation; per-element
operations are too slow due the additional interactive layer between
the operator and the physical machine.

And finally, there is a question of choice: Do we want to save some memory,
to be fast, or to be easy to understand?


\subsection{The effective QR factorisation}

This section describes algorithms for the effective QR factorisation
as has been linked above. The algorithms are designed for a heavy load
processing.

The effective QR factorisation in Algorithm~\ref{alg:qrfac} is
a transcription of Algorithm 5.2.1 by \cite{golub2013matrix}.
The function \textsf{housevec} is defined in Algorithm~\ref{alg:housevec}.
The function \textsf{housemul} is a convenient implementation
of (\ref{eq:golub514a}).
Having Householder vectors $v$ stored inside the factored matrix,
we can easy prepare the orthogonal matrix $Q$ at any moment;
Algorithm~\ref{alg:qform} (the equation (5.1.5) due \cite{golub2013matrix})
is such representation.

\begin{algorithm}
  \caption{The effective QR factorisation}
  \label{alg:qrfac}
  \begin{algorithmic}
    \Function{qrfac}{$A,\beta$}
    \State \Comment{Let $A \in \mathcal{R}^{m \times n}, \beta \in \mathcal{R}^n$}
    \State \Comment{$A$ contains a matrix to factorise; the upper
    triangle of $R$ and Householder vectors $v$ on output.
    $\beta$ is an output vector storing computed normalisation.}
    \State \Comment{An auxiliary variable: $v \in \mathcal{R}^m$}
    \For {$\; j = 1 \ldots n$}
    \State $v(\mdx{j}{m}), \beta_j$ = \Call{housevec}{$A(\mdx{j}{m},j)$}
    \State $A(\mdx{j}{m},\mdx{j}{n})$ = \Call{housemul}{$v(\mdx{j}{m}), \beta_j,
      A(\mdx{j}{m},\mdx{j}{n})$}
    \If {$\; j < m$}
    \State $A(\mdx{j+1}{m},j) = v(\mdx{j+1}{m})$
    \EndIf
    \EndFor
    \State Returns factored $A$ and $\beta$, see section.~\ref{subsec:memshare}.
    \EndFunction
  \end{algorithmic}
\end{algorithm}


\begin{algorithm}
  \caption{Householder matrix multiplication}
  \label{alg:housemul}
  \begin{algorithmic}
    \Function{housemul}{$v,\beta,A$}
    \State \Comment{Let $A \in \mathcal{R}^{m \times n}, v \in \mathcal{R}^n,
      \beta \in \mathcal{R}$}
    \State \Comment{Auxiliary variables: $u,w \in \mathcal{R}^m$}
    \State $u(\mdx{j}{m}) = \beta v(\mdx{j}{m})$
    \For {$\; j = 1 \ldots n$}
    \State $w_j = v^T \cdot A_j$
    \EndFor
    \For {$\; i = 1 \ldots m, j = 1 \ldots m$}
    \State $A_{ij} = A_{ij} - u_j w_i$
    \EndFor
    \State Returns $RP$ stored on place.
    \EndFunction
  \end{algorithmic}
\end{algorithm}


\begin{algorithm}
  \caption{Evaluate $Q$ from $v$ stored under diagonal of $R$}
  \label{alg:qform}
  \begin{algorithmic}
    \Function{qform}{$A,\beta$}
    \State \Comment{Let $A \in \mathcal{R}^{m \times n}, \beta \in \mathcal{R},
      Q \in \mathcal{R}^{m \times m}$}
    \State \Comment{Auxiliary variables: $v,u,w \in \mathcal{R}^m$}
    \State $Q = I_{nn}$
    \For {$\; j = n, n-1, \ldots, 1$}
    \State $v(\mdx{j}{m}) = (1, A(\mdx{j+1}{m},j))^T$
    \State $u(\mdx{j}{m}) = \beta_j v(\mdx{j}{m})$
    \For {$\; l = j \ldots n$}
    \State $w_l = v(\mdx{j}{m})^T \cdot Q(\mdx{j}{m},l)$
    \EndFor
    \For {$\; k = j \ldots m, l = j \ldots n$}
    \State $Q_{kl} = Q_{kl} - w_l u_k$
    \EndFor
    \EndFor
    \State Returns $Q$
    \EndFunction
  \end{algorithmic}
\end{algorithm}

All the presented algorithms are only slightly changed prototypes of
the equivalents of \cite{golub2013matrix} intentionally
keeping notations of the book.
Only a few important differences are here:
Householder vector $v$ is starting from $j$-th index, rather than
from the first index, to be consistent in notation.
The difference is observable on last line of
Algorithm~\ref{alg:qrfac} where $A(\mdx{j+1}{m},j) = v(\mdx{j+1}{m})$
is replaced by $A(\mdx{j+1}{m},j) = v(\mdx{1}{m-j+1}$).
Another difference is to keep of the vector $\beta$ for later use,
rather than recovering it during subsequent applications
to save computation time.

Also, all the
algorithms are arranged such way, to by easy implemented
in Fortran (see \texttt{src/qrhouse.f08}). Just for comparison,
the particular implementation of Algorithm~\ref{alg:housemul}
in Fortran can be found in listing~\ref{src:housemul}.

\begin{lstlisting}[float,language=Fortran,basicstyle=\ttfamily,
    caption={Householder matrix multiplication in Fortran},label={src:housemul}]

subroutine housemul(v,beta,A)

  real, dimension(:), intent(in) :: v
  real, intent(in) :: beta
  real, dimension(:,:), intent(in out) :: A

  real, dimension(size(v)) :: u, w
  integer :: i,j,m,n

  m = size(A,1)
  n = size(A,2)
  u = beta*v
  forall(j = 1:m) w(j) = dot_product(v,A(:,j))
  forall(i = 1:n, j = 1:m) A(i,j) = A(i,j) - u(i)*w(j)

end subroutine housemul
\end{lstlisting}




\section{Householder reflection}
\label{sec:house}

To nullify the desired elements of $R$ during $j$-th iteration
by the numerically stable way,
Householder vector $v^{(j)}$ is defined via the corresponding
column of $A$ as $x^{(j)} = A_j$ ($j$-th column of $A$)
\begin{equation} \label{eq:housevec}
  v_i^{(j)} =
  \begin{cases}
    0, & i < j,\\
    v_j, & i = j,\\
    x_j, & i > j,
  \end{cases}
\end{equation}
where
\begin{equation} \label{eq:housevec0}
  v_j = x_j \pm \sign(x_j) \lVert x(\mdx{j}{m}) \rVert_2.
\end{equation}
The sign in (\ref{eq:housevec0}) as well as the computation of $v_j$
are source of a disputation. \cite{golub2013matrix}
proposes negative sign (to maximise norm) which is keeping the numerical
precision. If $x$ is close to a positive multiple of $e$,
we can improve the precision by the way
\[
v = x - \lVert x \rVert_2 =
\frac{x_1^2 - \lVert x\rVert_2}{x_1 + \lVert x \rVert_2} =
-\frac{x_2^2 + \ldots + x_m^2}{x_1 + \lVert x \rVert_2}.
\]
Finally, we define
%The unit vector $e_k = (0 \ldots, 1, 0 \ldots 0)^T$.
\begin{equation} \label{eq:vj}
v_j =
\begin{dcases}
  -\frac{x_{j+1}^2 + \ldots + x_m^2}{x_j + \lVert x(\mdx{j}{m}) \rVert_2}, & x_j \ge 0,\\
  x_j - \lVert x(\mdx{j}{m}) \rVert_2, & x_j < 0.
\end{dcases}
\end{equation}
In the important case of a rank deficient matrix,
Householder vector will $v = 0$,
and we should add the condition $\beta = 0$.

\begin{algorithm}
  \caption{Householder vector}
  \label{alg:housevec}
  \begin{algorithmic}

    \Function{housevec}{$x$}
    \State \Comment{$x, v \in \mathcal{R}^m, \beta \in \mathcal{R}$}
    \State $v(\mdx{2}{m}) = x(\mdx{2}{m})$
    \State $\sigma^2 = x(\mdx{2}{m})^T \cdot x(\mdx{2}{m})$
    \If{$\, \sigma^2 > 0$}
    \State $d = \sqrt{x_1^2 + \sigma^2}$
    \If{$\, x_1 > 0$}
    \State $v_1 = -\sigma^2 / (x_1 + d)$
    \Else
    \State $v_1 = x_1 - d$
    \EndIf
    \If{$\, \lvert v_1 \rvert > 0$}
    \State $v = v / v_1$
    \State $\beta = 2 / (v^T \cdot v)$
    \Else
    \State $\beta = 0, v = 0$
    \EndIf
    \Else
    \State $v_1 = 1$
    \If{$\, x_1 \ge 0$}
    \State $\beta = 0$
    \Else
    \State $\beta = 2$
    \EndIf
    \EndIf
    \State \Return $v, \beta$
    \EndFunction
  \end{algorithmic}
\end{algorithm}


\begin{example}{Householder vector}
  \label{ex:housevec}
  We will compute Householder vector $v$ and matrix $P$
  for the first column of the matrix
  (example\footfullcite{wiki:qr} by wiki ,
  our procedure reflects the distinct definition):
  \begin{equation}
    \begin{pmatrix*}[r]
      12, & -51, & 4,\\
      6, & 167, & -68,\\
      -4, & 24, & -41.
    \end{pmatrix*}
  \end{equation}

  For $j = 1$, we have
  $v^{(1)} = (1, -3, 2)^T, \beta = \sfrac{1}{7}$
  from (\ref{eq:housevec}), (\ref{eq:vj}) and (\ref{eq:housemat})
  \begin{equation}
    H^{(1)} =
    \begin{pmatrix*}[r]
      1, & 0, & 0,\\
      0, & 1, & 0,\\
      0, & 0, & 1.
    \end{pmatrix*}
    - \frac{1}{7}
    \begin{pmatrix*}[r]
      1, & -3, &  2,\\
     -3, &  9, & -6,\\
      2, & -6, &  4.
    \end{pmatrix*} =
    \begin{pmatrix*}[r]
      \sfrac{6}{7},  & \sfrac{3}{7},  &  \sfrac{-2}{7},\\
      \sfrac{3}{7},  & \sfrac{-2}{7}, &  \sfrac{6}{7},\\
      \sfrac{-2}{7}, & \sfrac{6}{7},  &  \sfrac{3}{7}.
    \end{pmatrix*}
  \end{equation}

    (to be continued).
\end{example}


\section{QR factorisation with pivoting}

QR factorisation including pivoting is important for the numerical
stability, and significantly extends field of usage of the developed methods.
It permits to find a solution in such real cases when a system of linear
equations is badly conditioned or singular. It also suppress
accumulation of numerical rounding errors emerged from long summations
and products, the most frequent operations executed on matrices.

The implemented pivoting uses the same principles as the
column pivoting known for Gauss elimination technique. The pivots
are represented by a matrix $\Pi$ which collects permutations of rows of $A$.
$\Pi$ has all elements as zeros except the ones on the appropriate places
\begin{equation}
  A \Pi = Q R.
\end{equation}
A common used storage of the permutation matrix is a vector
$p \in \mathcal{N}_0^m$.
It is implemented in Algorithm~\ref{alg:qrpivo} which is
a rich variant of Algorithm~\ref{alg:qrfac}.
The permutation matrix $\Pi$ can be easy recovered from $p$ as
Algorithm~\ref{alg:perm} demonstrates.

\begin{algorithm}
  \caption{QR factorisation with column pivoting}
  \label{alg:qrpivo}
  \begin{algorithmic}
    \Function{qrpivot}{$A,\beta, p$}
    \State \Comment{Let $A \in \mathcal{R}^{m \times n}, \beta \in \mathcal{R}^n,
    p\in \mathcal{N}_0^m, m \ge n$.}
    \State \Comment{$A$ contains a matrix to factorise; the upper
    triangle of $R$ and Householder vectors $v$ and factors $\beta$ on output.
    $p$ is vector of permutations.
    }
    \State \Comment{An auxiliary variables: $c,v \in \mathcal{R}^m,
      \tau \in \mathcal{R}$}
    \State $p = (\{1, \ldots, n \} )^T$
    \For{$\; j =  1 \ldots n$}
    \State $c(j) = A(\mdx{1}{m},j)^T \cdot A(\mdx{1}{m},j)$
    \EndFor
    \For{$\, r = 1, \ldots, n$}
    \State Find the smallest index $r \le k \le n$ for which $c(k)$
    is the maximum.
    \If{$\, \mathrm{not} \, (c(k) > 0)$}
    \State exit loop over $r$
    \EndIf
    \If{$r \ne k$}
    \State $p(r) \leftrightarrow p(k)$
    \State $c(r) \leftrightarrow c(k)$
    \State $A(\mdx{1}{m},r) \leftrightarrow A(\mdx{1}{m},k)$
    \EndIf
    \State $v(\mdx{r}{m}), \beta(r)$ = \Call{housevec}{$A(\mdx{r}{m},j)$}
    \State $A(\mdx{r}{m},\mdx{r}{n})$ = \Call{housemul}{$v(\mdx{r}{m}), \beta(r), A(\mdx{r}{m},\mdx{r}{n})$}
    \State $A(\mdx{r+1}{m},r) = v(\mdx{r+1}{m})$
    \For{$\; l =  r+1 \ldots n$}
    \State $c(l) = c(l) - A(r,l)^2$
    \EndFor
    \EndFor
    \State \Return $A$ factored, $\beta,p$
    \EndFunction
  \end{algorithmic}
\end{algorithm}


\begin{algorithm}
  \caption{Compilation of the permutation matrix}
  \label{alg:perm}
  \begin{algorithmic}
    \State $\Pi = 0$
    \For{$\; j = 1, \ldots, m$}
    \State $\Pi(p(j),j) = 1$
    \EndFor
  \end{algorithmic}
\end{algorithm}



\section{A solution of a system of linear equations}
\label{sec:lineq}



A solution $x \in \mathcal{R}^m$ of a system of linear equations
\begin{equation}
  Ax = b
\end{equation}
with a matrix $A$ factored on $Q,R$, and a right side $x \in \mathcal{R}^m$,
is straightforward thanks to the identity
\begin{equation}
  Q^{-1} = Q^T.
\end{equation}
The full system $A$ is transformed on the system having triangle matrix $R$
\begin{equation}
  Rx = q,
\end{equation}
and the right side $q = Q^T b$. To solve one,
we can apply the back-substitution under the condition
$\lvert R_{ii} \rvert \ne 0$ (see Algorithm~\ref{alg:rsol}):
\begin{equation}
  x_i = \frac{1}{R_{ii}} \left(
  q_i - \sum_{j=i+1}^m R_{ij} x_j \right), \qquad i = m, \ldots, 1.
\end{equation}

\begin{algorithm}
  \caption{Solution of $Rx=q$ with back-substitution}
  \label{alg:rsol}
  \begin{algorithmic}
    \Function{rsol}{$A,q$}
    \State \Comment{Let $A \in \mathcal{R}^{m \times n}, q,u \in \mathcal{R}^n$}
    \For{$\, i = m, \ldots, 1$}
    \State $u(i) = [q(i) - A(i,\mdx{i+1}{n})^T \cdot x(\mdx{i+1}{n})] / A(i,i)$
    \EndFor
    \State \Return $u$
    \EndFunction
  \end{algorithmic}
\end{algorithm}


Note that $Q$ can be again restored from vectors $v$ stored in columns of $R$
by applications of the trick presented in section~\ref{sec:speedup}.
Algorithm~\ref{alg:houseapp} materialises the idea.

This way is only valid for full square matrices with dimensions
$m \times m$. This is, in fact, an equivalent to direct Gaussian elimination
or its variants as $LU$ or $LL^T$ factorisation, notwithstanding
a numerical class of the matrices is wider (\cite{golub2013matrix}, Sect.~5.3).

The essential advantage of QR factorisation is revealed for
over-determined systems having more equations than variables ($m > n$),
and under-determined systems having less equations than variables ($m < n$).
The particular important case are rank deficient singular matrices.

Next paragraphs briefly describes solutions in the difficult cases,
which covers Algorithm~\ref{alg:qrsol}.


\subsection{A least square solution of a system of linear equations}
\label{sec:lsqrt}

The least square solution minimises of Pythagorean norm,
a generalised distance, between a computed $Ax$, and expected $b$,
solutions
\begin{equation}
  \min \lVert A x - b \rVert_2,
\end{equation}
where $A \in \mathcal{R}^{m \times n}, m \ge n, b \in \mathcal{R}^m,
x \in \mathcal{R}^n$.
The problem can be solved with help of QR factorisation which decompose
the system on the two parts:
\begin{equation}
  Q^T A = R =
  \begin{pmatrix*}[c]
    R_1,\\
    0
  \end{pmatrix*},
  \quad Q^T b =
  \begin{pmatrix*}[c]
    c,\\
    d
  \end{pmatrix*}.
\end{equation}
$R_1 \in \mathcal{R}^{n \times n}$ is a square matrix. The lower part
of factorised $R$ is nullified. $c \in \mathcal{R}^n$ part corresponds
to the solution, while $d \in \mathcal{R}^{m-n}$ stores residuals.
\begin{equation}
  \lVert A x - b \rVert_2^2 = \lVert Q^T A x - Q^T b \rVert_2^2 =
  \lVert R_1 x - c \rVert_2^2 + \lVert d \rVert_2^2.
\end{equation}
The solution in sense of least squares is than
\begin{equation}
  R_1 x_0 = c,
\end{equation}
and the residual sum is
\begin{equation}
\lVert Ax_0 - b \rVert_2 = \lVert d \rVert_2.
\end{equation}



\begin{algorithm}
  \caption{Householder matrix application}
  \label{alg:houseapp}
  \begin{algorithmic}
    \Function{houseapp}{$A,\beta,x$}
    \State \Comment{Let $A \in \mathcal{R}^{m \times n}, v,\beta \in \mathcal{R}^n$}
    \For{$\, j = 1, \ldots, m$}
    \State $v(\mdx{j}{m}) = (1, A(\mdx{j+1}{m},j))^T$
    \State $x(\mdx{j}{m}) = x(\mdx{j}{m}) - \beta_j
           [v(\mdx{j}{m})^T \cdot x(\mdx{j}{m})] v(\mdx{j}{m})$
    \EndFor
    \State \Return $x$
    \EndFunction
  \end{algorithmic}
\end{algorithm}


\begin{algorithm}
  \caption{Linear equations solver}
  \label{alg:qrsol}
  \begin{algorithmic}
    \Function{qrsol}{$A,b,x$}
    \State \Comment{Let $A,R \in \mathcal{R}^{m \times n}, x,b \in \mathcal{R}^n$}
    \State \Comment{Auxiliary variables: $u,q,\beta \in \mathcal{R}^m$}

    \If{$m > n$}
    \State \Comment{The least-square solution, Section~\ref{sec:lsqrt}}
    \State $R,\beta = $ \Call{qrfac}{$A$}
    \State $q =$  \Call{houseapp}{$R,\beta,b$}
    \State \Return $x =$  \Call{rsol}{$R(\mdx{1}{n},\mdx{1}{n}),q$}
    \Else
    \State $R,\beta,\Pi = $ \Call{qrpivot}{$A$}
    \State $q = Q^T b = $  \Call{houseapp}{$R,\beta,b$}

    \State Set $\delta > 0, \delta = \tau \varepsilon \lvert R(1,1) \rvert,
    R(1,1) \ne 0$
    \State Determine rank: $r = \{ i, \lvert R(i,i) \rvert > \Delta \}$
    \If{$r = m$}
    \State \Comment{A square matrix Section~\ref{sec:lineq}}
    \State $z =$ \Call{rsol}{$R,q$}
    \State \Return $x = \Pi^T z$
%    \State $u =$  \Call{rsol}{$R,q$}
%    \State \Return $x = u(p)$
    \ElsIf{$r < m$}
    \State \Comment{A rank deficient matrix, Section~\ref{subsec:lsrank}}
    \State \Comment{Basic solution: $u(\mdx{1}{r}) =$
      \Call{rsol}{$R(\mdx{1}{r},\mdx{1}{r}),q(\mdx{1}{r})$}, $x_B = \Pi^T u$}
    \State $S(\mdx{1}{n},\mdx{1}{r}),\beta'(\mdx{1}{r}) = $
    \Call{qrfac}{$R^T(\mdx{1}{r},\mdx{1}{n})$}
    \State $U(\mdx{1}{n},\mdx{1}{n})=$ \Call{qform}{$S(\mdx{1}{n},\mdx{1}{r}),\beta'(\mdx{1}{r})$}
    \State $z(\mdx{1}{r}) = $ \Call{lsol}{$S(\mdx{1}{r},\mdx{1}{m}),q(\mdx{1}{r})$}
    \State \Return $x = \Pi^T U(:,\mdx{1}{r}) z(\mdx{1}{r}) $
    \EndIf
    \EndIf

    \EndFunction
  \end{algorithmic}
\end{algorithm}


\subsection{A basic solution of linear system with a rank deficient matrix}
\label{subsec:basic}

We know that rank deficient matrix has no unique solution.
The class of the possible solutions is usually described
with help of a slack variable, known as a parameter,
for analytical solutions.

The parametric solution approach is not suitable for a numerical
solution. It is reason than one adds an additional information
to select some particular solution. A basic solution is developed
in this section, while next section describes a solution
with the minimal norm.

A solution of rank deficient problems has started by column pivoted
QR factorisation $A \Pi = QR, Q^{-1} = Q^T$. The pivoting process sorts columns
by size of diagonal elements
$\lvert R_{11} \rvert \ge \lvert R_{22} \rvert \ge \ldots \ge \lvert R_{mm} \rvert $.
Potentially linearly dependent columns has the diagonal elements close
to the machine epsilon $\varepsilon$. In such case, a condition
($\tau > 1$), says
\begin{equation}
  \lvert R_{ii} \rvert < \tau \varepsilon \lvert R_{11} \rvert
\end{equation}
can be used to reveal a numerical rank of the matrix. We will
denote the rank, an index of a last non-zero diagonal element in numerical
sense, by $r$,

Let $r \in \mathcal{N}$ is the rank of matrix $A$,
the last index for which
$\lvert R_{rr} \rvert > \varepsilon \lvert R_{11} \rvert$.
The pivoting process will split the matrix $R$, and the system, on two parts
\begin{equation}
  A \Pi = Q
  \begin{pmatrix*}[c]
    R_{11}, & R_{12}, \\
    0 & 0
  \end{pmatrix*}.
\end{equation}
with $R_{11} \in \mathcal{R}^{r \times r}, R_{12} \in \mathcal{R}^{r \times m-r}$.
The least square approach requires
\begin{equation}
  \lVert Ax - b\rVert_2^2 = \lVert (Q^T A \Pi) (\Pi x) - (Q^Tb) \rVert_2^2
  = \lVert R_{11} y (c - R_{12} z) \rVert_2^2 + \lVert  d \rVert_2^2,
\end{equation}
where
\begin{equation}
  \Pi^T x =
  \begin{pmatrix*}[c]
    y,\\
    z
  \end{pmatrix*},
  \qquad
  Q^T b =
  \begin{pmatrix*}[c]
    c,\\
    d
  \end{pmatrix*}.
\end{equation}
The solution is
\begin{equation}
  x_B = \Pi
  \begin{pmatrix*}[c]
    R_{11}^{-1} (c - R_{12} z),\\
    z
  \end{pmatrix*}.
\end{equation}
Ones calls $x_\mathrm{B}$ solution as a basic solution. The basic solution
is no solution in least square sense, except case $R_{12}=0$.


\subsection{A solution in least square sense by a complete orthogonal decomposition}
\label{subsec:lsrank}

The least square solution of rank deficient matrix
can be found by application of a dark magic trick.
If we continue with the revealed information from previous
paragraph \ref{subsec:basic}, the rectangular matrix
has dimensions $r \times m$. We can transform the matrix on the least
square solution studied in Section~\ref{sec:lsqrt} by a matrix transposition
having dimensions $m \times r, r \le m$.

By the way, we apply again the QR factorisation on already factored
transposed triangle of $R$: $R^T = US$
\begin{equation}
  \begin{pmatrix*}[c]
    R_{11}^T,\\
    R_{12}^T
  \end{pmatrix*}
  = U
  \begin{pmatrix*}[c]
    S_1,\\
    0
  \end{pmatrix*}
\end{equation}
and so the system $A \Pi  x=b$ can be rewritten on $R \Pi x = Q^T b$,
and than on
\begin{equation}
  S_1 z = Q^T b
\end{equation}
and so
\begin{equation}
  \Pi x_{\mathrm{LS}} = U z
\end{equation}
Finally, the permutation should be inverted $x = u(p).$

This way factorises a matrix on two orthogonal spaces described
by $U,V$ matrices, or more exactly
\begin{equation}
  A = U
  \begin{pmatrix*}[c]
    T_{11}, & 0\\
    0, & 0
  \end{pmatrix*}
  V^{-1}
\end{equation}
where $V = \Pi Q$, and we are identifying $T_{11} = S_1^T$.
The process is known as the complete orthogonal decomposition.

This framework gives exactly the same solution as well-known SVD,
by significantly cheaper way: The number of operations required
for a single QR run is $~2 n^3/3$ compared to $~12n^3$
(\cite{golub2013matrix}, p.~293). The drawback of SVD
is the interactive algorithm and more deep matrix analysis ---
singular values (eigenvalues) and eigenvectors are hard to find.
By the way, a natural choice for a quick solution of systems with
a (potentially) rank-deficient matrix is the described complete
orthogonal decomposition.


\begin{algorithm}
  \caption{Solution of $Lx=q$ with back-substitution}
  \label{alg:lsol}
  \begin{algorithmic}
    \Function{lsol}{$L,q$}
    \State \Comment{Let $L \in \mathcal{R}^{m \times n}, q,x \in \mathcal{R}^n$}

    \For{$\, i = 1, \ldots, m$}
    \State $x(i) = [q(i) - L(\mdx{1}{i-1},i)^T \cdot x(\mdx{1}{i-1})] / L(i,i)$
    \EndFor
    \State \Return $x$
    \EndFunction
  \end{algorithmic}
\end{algorithm}

\begin{algorithm}
  \caption{Determine rank of a matrix}
  \label{alg:rank}
  \begin{algorithmic}
    \Function{rank}{$R,\tau$}
    \State \Comment{Let $R \in \mathcal{R}^{m \times n}, \tau > 0 \in \mathcal{R}$}
    \State $r = 0$
    \For{$\, i = 1, \ldots, m$}
    \If{$\, \lvert R(i,i) \rvert < \tau$}
    \State \Return $r$
    \EndIf
    \State $r = i$
    \EndFor
    \State \Return $r$
    \EndFunction
  \end{algorithmic}
\end{algorithm}

\section{About a matrix inversion}

Let $A \in \mathcal{R}^{m \times m}$ is a square matrix. The inverse matrix
satisfies these conditions
\begin{equation}
  \label{eq:invcond}
  A A^{-1} = I, \quad A^{-1} A = I.
\end{equation}

Let $A$ is a full rank matrix. A way for computing of inverse matrix
is a solution of series of linear equations with unit vectors.
The unitary matrix $I$ has unit vectors $e_j$ in columns (rows).
By the way, for $j$-th column of $a_j = A_j^{-1}$, we have this linear
system
\[
A a_j = e_j, \qquad j=1, \ldots, n.
\]
which can be solved for $n$ times to fill $a_j = A_j^{-1}$.
This is just a direct generalisation of methods mentioned in
previous section~\ref{sec:lineq}.

Note that solving systems for coordinate vectors $e_i$ removes the need for
the matrix multiplication of right side:
all elements $e_j$ are zeros except the one in $j$-th row:
\[
Q^T e_j = Q^T(:,j),
\]
so the operation just selects the appropriate column of $Q$

The rank deficient problems, or problems with non-square matrix,
can be maintained by similar way as we study in the solutions
of Sect.~\ref{subsec:lsrank}.
The result will no more satisfy conditions (\ref{eq:invcond});
they are replaced by the conditions:
\[
A A^+ A = A, \ A^+ A A^+ = A^+, \  (A A^+)^T = A A^+, \ (A^+ A)^T = A^+ A.
\]
Such matrix is known as pseudo-inverse matrix
$A^+$ having minimal Frobenius norm
\[
\lVert A A^+ - I \rVert_F.
\]
The pseudo-inverse is very suitable to estimate of
a covariance (sub-)matrix.
The final algorithm~\ref{alg:qrinv} works for both full-rank
and rank deficient systems.

\begin{algorithm}
  \caption{(Pseudo-)Inverse of a matrix}
  \label{alg:qrinv}
  \begin{algorithmic}
    \Function{qrinv}{$A,A^{-1},\tau$}
    \State \Comment{Let $A,A^{-1},A^+ \in \mathcal{R}^{m \times n},
      x,b \in \mathcal{R}^m, \tau > 1 \in \mathcal{R}$}
    \State \Comment{Auxiliary variables: $Q,U,W \in \mathcal{R}^{m \times n},
      \Pi \in \mathcal{N}^{m \times m},
      z,q,\beta,\beta' \in \mathcal{R}^m$}

    \State $R,\beta,\Pi = $ \Call{qrpivot}{$A$}
    \State $Q = $ \Call{qform}{$R,\beta$}
    \State $r = $ \Call{rank}{$A, \tau \varepsilon \lvert R(1,1) \rvert$}

    \If{$\, r =  m$}
    \For{$\, i = 1, \ldots, m$}
    \State $z =$ \Call{rsol}{$R,Q(i,:)$}
    \State $A^{-1}(:,i) = \Pi^T z$
    \EndFor
    \State \Return $A^{-1}$
    \ElsIf{\, $r < m$}
    \State $S(\mdx{1}{n},\mdx{1}{r}),\beta'(\mdx{1}{r}) = $
    \Call{qrfac}{$R(\mdx{1}{r},\mdx{1}{n})^T$}
    \State $U(\mdx{1}{n},\mdx{1}{n})=$
    \Call{qform}{$S(\mdx{1}{n},\mdx{1}{r}),\beta'(\mdx{1}{r})$}
    \For{$\, i = 1, \ldots, n$}
    \State $z(\mdx{1}{r}) = $ \Call{lsol}{$S(\mdx{1}{r},\mdx{1}{r}),Q(i,\mdx{1}{r})$}
    \State $A^+(:,i) = \Pi^T U(:,\mdx{1}{n}) z$
    \EndFor
    \State \Return $A^+$
    \EndIf
    \EndFunction
  \end{algorithmic}
\end{algorithm}



\section{Eigenvalues and eigenvectors of a symmetric matrix}

The framework of QR factorisation can be utilised also
on a problem of determination of eigenvalues and
eigenvectors of a symmetric square matrix $A \in \mathcal{R}^{n \times n}$.
Eigenvalues $\lambda$ are roots of a characteristics
polynomial taken from determinant of a matrix
\begin{equation}
  \lvert A  - \lambda I \rvert = 0.
\end{equation}
Counts of the roots is exactly the order of the polynomial $n$.
The roots are real numbers $\lambda_i \in \mathcal{R}$ for
symmetric real matrices.

An eigenvector $z_i \in \mathcal{R}^n$ is a vector, which is associated
to every eigenvalue via a solution of the system of equations
\begin{equation} \label{eg:eigenvec}
  A z_i = \lambda_i z_i, \quad i = 1, \ldots n.
\end{equation}
The eigenvectors are orthogonal each other $z_i \perp z_j, z_i^T \cdot z_j = 0$.

The characteristic polynomial has analytic solution only
up to order three. All others roots should be, in principle,
found by an iterative technique. That is the reason, why I am
limiting only on symmetric problems having efficient methods developed;
fortunately, only such kind of problems is encountered in field of optimisation.

The basic Algorithm~\ref{alg:power} for determination of eigenvalues is
on base of the power method (cite?): reiteration of QR factorisation.
Eigenvalues are then computed by direct solution of (\ref{eg:eigenvec});
the system is singular by definition, so I use the parametric solution
for the demonstration.


\begin{algorithm}
  \caption{Outline for determination of eigenvalues and vectors}
  \label{alg:power}
  \begin{algorithmic}
    \State \Comment{Let $A,Q,R,\Lambda,Z \in \mathcal{R}^{m \times n}, b \in \mathcal{R}^n$}
    \State $\Lambda = A$
    \For{$\; k = 1, \ldots$}
    \State $Q,R = $ \Call{qrfac}{$\Lambda$}
    \State $\Lambda = Q R$
    \State $d = \max\{\lvert \Lambda_{\div} \rvert\}$,
    $\Lambda_{\div}$ denotes off-diagonal elements
    \State Finish when $d$ is under a tolerance
    \EndFor
    \For{$\; i = 1, \ldots, n$}
    \State $Q,R = $ \Call{qrfac}{$A - \lambda_i I$}
    \State $Z_{ni} = \lVert \{ R_{ii}, i = 1, \ldots, n \} \rVert_2$
    \For{$\; j = n-1, \ldots, 1$}
    \State $Z_{ji} = -[R(j,\mdx{j+1}{n})^T \cdot Z(j,\mdx{j+1}{n})] / R_{jj}$
    \EndFor
    \State $Z_i = Z_i / \lVert Z_i \rVert_2$, for a non-zero vector
    \EndFor
    \State $\lambda_i = \{ \Lambda_{ii}, i = 1, \ldots, n \}$
    \State \Return $\lambda, Z$
  \end{algorithmic}
\end{algorithm}

The drawbacks of the algorithm ~\ref{alg:power} are:
the slow convergence, and the needs for additional eigenvectors computations.
There is again an effective way how to compute both eigenvalues and vectors,
like in case of QR factorisation.
The algorithm, described in this section,
computes both values and vectors simultaneously
within a minimal number of iterations.


\section[An effective eigenvalue decomposition]{An effective algorithm for
  determination of eigenvalues and eigenvectors of a symmetric matrix on base
  of QR factorisation}

Today convenient algorithms are on base of application of elementary
rotations which nullifies off-diagonal elements: if the elemental rotation
are applied repeatedly, they effective transforms an original matrix
onto diagonal form:
\begin{equation}
  D = Z^T A Z.
\end{equation}
The matrices $Z$ contains eigenvectors, while eigenvalues are
on the diagonal. As for QR factorisation, it is possible to
save a computer memory and store appropriate matrices in which
are overwrote.

%transformace, ukazka vypoctu definice parametru



The elementary rotations --- known as Givens transformation ---
are the matrices:
\begin{equation}
  \begin{pmatrix*}[c]
    1,     & \cdots & 0,     & \cdots & 0,     & \cdots & 0,\\
    \vdots & \ddots & \vdots & {}     & \vdots & {}     & \vdots\\
    0,     & \cdots & c,     & \cdots & s,     & \cdots & 0,\\
    \vdots & {}     & \vdots & \ddots & \vdots & {}     & \vdots\\
    0,     & \cdots & -s,    & \cdots & c,     & \cdots & 0,\\
    \vdots & {}     & \vdots & {}     &\vdots  & \ddots & \vdots\\
    0,     & \cdots & 0,     & \cdots & 0,     & \cdots & 1\\
  \end{pmatrix*}.
\end{equation}
If we know the coordinates $x,y$, the rotation, the angle in $-\pi \ldots \pi$
can be determined as Algorithm~\ref{alg:givens} summarises.

\begin{algorithm}
  \caption{Givens rotation}
  \label{alg:givens}
  \begin{algorithmic}
    \Function{givens}{$x,y$}
    \If{$\; \lvert y \rvert < \varepsilon$}
    \State $c = 1, s = 0$
    \Else
    \If{$\; \lvert y \rvert > \lvert x \rvert$}
    \State $\tau = -x/y, s = 1 / \sqrt{1 + \tau^2}, c = s \tau$
    \Else
    \State $\tau = -y/x, c = 1 / \sqrt{1 + \tau^2}, s = c \tau$
    \EndIf
    \EndIf
    \State \Return $c,s$
    \EndFunction
  \end{algorithmic}
\end{algorithm}





\subsection{Householder tri-diagonalisation}
\label{subsec:housetridig}

Every symmetric matrix can be transformed by application of Householder
rotations onto tri-diagonal matrix. The transformation applies equivalent
matrices which does not changes eigenvalues. If we are able to collect
the transformations, we can determine eigenvectors. The tri-diagonalisation
significantly simplifies consecutive iteration process via QR factorisation.


\begin{algorithm}
  \caption{Householder tri-diagonalisation}
  \label{alg:housetridig}
  \begin{algorithmic}
    \Function{housetridig}{$A$}
    \State \Comment{Let $A,T \in \mathcal{R}^{n \times n}, a\in \mathcal{R}^n, b \in \mathcal{R}^{n-1},$}
    \State \Comment{An auxiliary variables: $p,v,w,u,\beta \in \mathcal{R}^n$}
    \State $T=A$
    \For{$\; k =  1, \ldots, n-2$}
    \State $v(\mdx{k+1}{n}), \beta_k$ = \Call{housevec}{$T(\mdx{k+1}{n},k)$}
    \State $p(\mdx{k+1}{n}) = \beta_k T(\mdx{k+1}{n},\mdx{k+1}{n}) \cdot v(\mdx{k+1}{n})$
    \State $r = \beta_k [p(\mdx{k+1}{n})^T \cdot v(\mdx{k+1}{n})] / 2$
    \State $w(\mdx{k+1}{n}) = p(\mdx{k+1}{n}) - r v(\mdx{k+1}{n})$
    \State $T_{k+1,k} = T_{k,k+1} = \lVert T(\mdx{k+1}{n},k) \rVert_2$
    \For{$\; i = k+1, \ldots, n, \, j = k+1, \ldots, n$}
    \State $T_{i,j} = T_{i,j} - v_i w_j - w_i v_j$
    \EndFor
    \If{$\; k > 1$}
    \State $T(\mdx{k+1}{n},k-1) = u(\mdx{k+1}{n})$
    \EndIf
    \State $u(\mdx{k+1}{n}) = v(\mdx{k+1}{n})$
    \EndFor

    \If{$\; n > 2$}
    \State $T_{n,n-2} = u_n$
    \EndIf
    \State $a = \{ T_{ii}, i = 1, \ldots, n \}$
    \State $b = \{ T_{i,i+1}, i = 1, \ldots, n-1 \}$
    \State $T(\mdx{n-1}{n},\mdx{n-1}{n}) = I$

    \For{$\; k = n-2, \ldots, 1$}
    \State $v(\mdx{k+1}{n}) = (1, T(\mdx{k+2}{n},k))^T$
    \State $T_{kk} = 1$
    \State $T(\mdx{k+1}{n}, k) = T(k,\mdx{k+1}{n}) = 0$
    \State $w(\mdx{k+1}{n}) = \{ [T(\mdx{k+1}{n},i)^T \cdot v(\mdx{k+1}{n})] , i = k+1, \ldots, n \}$
    \State $u(\mdx{k+1}{n}) = \beta_k v(\mdx{k+1}{n})$
    \For{$\; i = k+1, \ldots, n, \, j = k+1, \ldots, n$}
    \State $T_{i,j} = T_{i,j} - w_j u_i$
    \EndFor
    \EndFor
    \State Returns $a,b,T$
    \EndFunction
  \end{algorithmic}
\end{algorithm}

The result is a symmetric tri-diagonal matrix and a matrix storing
the transformations. The economy use of memory follows ideas developed
for QR factorisation. The tri-diagonal matrix is can be stored in
two vectors $a \in \mathcal{R}^n, b \in \mathcal{R}^{n-1}$:
\begin{equation}
  \begin{pmatrix*}[c]
    a_1, & b_1,   & 0,   & 0        & \ldots, & 0\\
    b_1, & a_2,   & b_2, & 0        & \ldots, & 0\\
    0,   & b_2,   & a_3, & b_3,     & \ldots, & 0\\
  \vdots & {}     & \ddots & \ddots & \ddots  & \vdots\\
    0,   & \dots, & 0,   & b_{n-2}  & a_{n-1}, & b_{n-1}\\
    0,   & \dots, & 0,   & 0       & b_{n-1}, & a_n\\
  \end{pmatrix*}.
\end{equation}



\subsection{Wilkinson shift}


\begin{algorithm}
  \caption{Implicit symmetric Francis QR step with Wilkinson shift}
  \label{alg:qrstep}
  \begin{algorithmic}
    \Function{qrstep}{$a,b,Q$}
    \State \Comment{Let $Q \in \mathcal{R}^{n \times n}, a\in \mathcal{R}^n, b \in \mathcal{R}^{n-1},$}

    \State $d = (a_{n-1} - a_n) / 2$
    \State $\mu = a_n - b_{n-1}^2 / [d + \sign(d) \lVert (d,b_{n-1}) \rVert_2]$
    \State $x = a_1 - \mu$
    \State $y = b_1$
    \For{$\; k = 1, \ldots, n-1$}
    \State $c,s =$ \Call{givens}{$x,y$}
    \State $u = a_k, v = a_{k+1}, w = 2 c s b_k$
    \State $a_k = c^2 u + s^2 v - w$
    \State $a_{k+1} = c^2 v + s^2 u + w$
    \State $b_k = (c^2 - s^2) b_k + cs(u-v)$
    \State $b_{k-1} = cx - sy$, if $k>1$
    \If{$\; k < n-1$}
    \State $x = b_k, y = -sb_{k+1},b_{k+1} = c b_{k+1}$
    \EndIf
    \State $Q(\mdx{1}{n}, \mdx{k}{k+1}) = Q(\mdx{1}{n}, \mdx{k}{k+1}) \begin{pmatrix*}[r] c, & s,\\ -s, & c \end{pmatrix*}$
    \EndFor
    \State Returns modified $a,b,Q$
    \EndFunction
  \end{algorithmic}
\end{algorithm}

Algorithm~\ref{alg:qrstep} is adapted for the tri-diagonal
matrix discussed in~\ref{subsec:housetridig}. The key operation
is the Givens rotation applied on the matrix:
\[
  \begin{pmatrix*}[c]
    c, & s \\
    s, & c \\
  \end{pmatrix*}
  \begin{pmatrix*}[c]
    T_{kk}, & T_{k,k+1} \\
    T_{k+1,k}, & T_{k+1,k+1} \\
  \end{pmatrix*}
  \begin{pmatrix*}[c]
    c, & s \\
    -s, & c \\
  \end{pmatrix*}.
\]
The application make a new non-zero element out of the tri-diagonal;
unfortunately, the immediate step will nullify then. The anomaly
is moved from up to down and will vanish after the last operation,
leaving clear tri-diagonal.

If the $a, b$ representation is used, the non-zero element should
by treated separately. The transformation
\begin{equation}
  \begin{pmatrix*}[c]
    1, & 0, &  0, & 0\\
    0, & c, & -s, & 0\\
    0, & s, &  c, & 0\\
    0, & 0, &  0, & 1
  \end{pmatrix*}
  \begin{pmatrix*}[c]
    a_{k-1}, & b_{k-1}, &  z, & 0\\
    b_{k-1}, & a_k, &    b_k, & 0\\
    z, & b_k, &  a_{k+1}, & b_{k+1}\\
    0, & 0, & b_{k+1}, & a_{k+1}
  \end{pmatrix*}
  \begin{pmatrix*}[c]
    1, & 0, &  0, & 0\\
    0, & c, &  s, & 0\\
    0, &-s, &  c, & 0\\
    0, & 0, &  0, & 1
  \end{pmatrix*}
\end{equation}
leads to the resultant matrix
\begin{equation}
  \begin{pmatrix*}[c]
    a_{k-1}, & cb_{k-1} - sz, & sb_{k-1} + zc, & 0\\
    cb_{k-1} - sz, & a_k', & b_k', & -sb_{k+1}\\
    sb_{k-1}+cz, & b_k', & a_{k+1}', & cb_{k+1}\\
    0, & -sb_{k+1}, & c b_{k+1}, & a_{k+2}
  \end{pmatrix*},
\end{equation}
where the updated elements of the tri-diagonal are
\begin{align}
  a_k' =& c^2 a_k + s^2 a_{k+1} - 2csb_k,\\
  a_{k+1}' =& c^2 a_{k+1} + s^2 a_k + 2csb_k,\\
  b_k' =& (c^2 - s^2) b_k + cs (a_k - a_{k+1}).
\end{align}
We also see that there are conditions
\begin{equation}
  b'_{k-1} = c b_{k-1} - s z, \qquad b'_{k+1} = c b_{k+1}.
\end{equation}



\subsection{A deflation}

A selecting of suitable sub-matrix during iterative process is know as deflation
or bounding. The off-diagonal elements of tri-diagonal matrix represented by
$b$ are nullified from the end or begin. An immediate matrix looks like
\begin{equation}
  \begin{pmatrix*}[c]
    a_1, & <\varepsilon,   & 0,   & 0        & \ldots, & 0\\
    <\varepsilon, & a_2,   & b_2, & 0        & \ldots, & 0\\
    0,   & b_2,   & a_3, & b_3,     & \ldots, & 0\\
    \vdots & {}     & \ddots & \ddots & \ddots  & \vdots\\

    0,   & \dots, & 0,   & <\varepsilon  & a_{n-1}, & <\varepsilon\\

    0,   & \dots, & 0,   & <\varepsilon  & a_{n-1}, & <\varepsilon\\
    0,   & \dots, & 0,   & 0       & <\varepsilon, & a_n\\
  \end{pmatrix*}.
\end{equation}
With removed the already diagonalised leading and tailing sub-matrices,
the computation speed can be improved, and numerical errors suppressed.



\begin{algorithm}
  \caption{Determines eigenvalues and eigenvectors.}
  \label{alg:qreig}
  \begin{algorithmic}
    \Function{qreig}{$A,\lambda,Z$}
    \State \Comment{Let $A,Q \in \mathcal{R}^{n \times n},
      \lambda,a\in \mathcal{R}^n, b \in \mathcal{R}^{n-1},$}
    \State $a,b,Q =$ \Call{housetridig}{$A$}
    \State $l = 1$
    \State $m = n$
    \For{$\; k = 1, \dots$}
    \For{$\; i = 2, \ldots, n-1$}
    \State $\delta(i) = \varepsilon(\lvert a_i \rvert + \lvert a_{i+1} \rvert)$
    \EndFor
    \For{$\; i = 2, \ldots, n-1$}
    \If{$\; \lvert b_i \rvert > \delta_i \; \text{and} \;
      \lvert b_{i-1} \rvert < \delta_{i-1}$}
    \State $l = i$
    \EndIf
    \If{$\; \lvert b_i \rvert < \delta_i \; \text{and} \;
      \lvert b_{i-1} \rvert > \delta_{i-1}$}
    \State $l = i$
    \EndIf
    \EndFor

    \If{$\; m > l$}
    \State \Call{qrstep}{$a(\mdx{l}{m}),b(\mdx{l}{m-1}),Q(\mdx{1}{n},\mdx{l}{m-1})$}
    \EndIf

    \State Finish, if $m - l = 1$
    \EndFor
    \State $\lambda = a$
    \State Returns modified $\lambda,Q$
    \EndFunction
  \end{algorithmic}
\end{algorithm}



%\section{John F. Francis}


\section{Bibliography notes}
\label{sec:qrbiblio}

The textbook of {} \textcite{ralston2012first} gives an introduction into
whole field of numerical mathematics including matrix decomposition itself.
\cite{golub2013matrix} is a comprehensive summary of general computation
methods of linear algebra; Lapack computer library\footfullcite{website:lapack}
is also based on the book.

There are many excellent on-line resources: blogs, documents or forums.
Mr. Higham\footfullcite{website:nhigham} provides the valuable collection
of related articles on his website.
Curses\footfullcite{website:arbenz} by Mr. Arbenz was been very helpful
for me during the implementation phase.


\printbibliography[heading=subbibliography]
\end{refsection}
